\section{The Grammar of First-Order Logic}

\subsection{Components of Syntax}

Last week, we explored this ins and outs of first-order logic informally. Now
we'll start to formalize those ideas, starting with the basic grammar of FOL.

\begin{definition}[Expressions of First-Order Logic]: We can difine an
  \label{def:fol-expression}
  \emph{expression} of FOL as any sequence of its six kinds of basic symbols:\\
  \begin{enumerate}
    \item \textbf{Predicates:} $A,B,C,\dots,=$
    \item \textbf{Names:} $a,b,c,\dots$
    \item \textbf{Variables:} $x,y,z,\dots$
    \item \textbf{Connectives:} $\enot,\eand,\eor,\eif,\eiff$ 
    \item \textbf{Quantifiers:} $\eforall,\eexists$
    \item \textbf{Brackets:} (,)
  \end{enumerate}
  The alphabetic symbols (names, predicates and variables) can be used with
  or without numeric subscripts, for example $A_1, A_2, A_3, \dots$ or
  $x_1, x_2, x_3, \dots$.
\end{definition}

We assume these basic expressions have roughly the meaning discussed informally.
But, of course, that is not important here. All we want to do with a grammar is
to say which expressions are sentences and which are not.

In the metalanguage we will refer to expressions with lower case greek letters.

The most basic basic expressions in FOL are \emph{singular terms}. Expressions
of this kind are noteworthy in their semantics: they each \emph{refer} to an
individual in the domain. So we will keep them distinct in the syntax.

\begin{definition}[Term]
  \label{def:fol-term}
  A \emph{term} in first-order logic is any \emph{name}. A name in FOL can be
  one of two kinds:
  \begin{enumerate}
  \item
    \textbf{Dummy Names:} lower case $a$ and $c$, with or without subscripts
    \textbf{Proper Names:} lower case latin letters, but not $a$ or $c$
  \end{enumerate}
\end{definition}

Dummy names are placeholders, like variables. But unlike variables, they are
full-fledged terms: we assume that they each refer directly to an individual
in the domain. Note that on this definition, variables are not terms. This is
as it should be, since they do not \emph{refer}. More on this when we get to
the semantics. For now, it is noteworthy that variables are used wherever
terms are allowed. They are not terms, but artefacts of the process that
introduces quantifiers. Before we talk about that process, let's work through
a few more preliminary definitions.

\begin{definition}[Predicate]
  An \emph{n-place predicate} in first-order logic is an expression which
  takes $n$ terms as \emph{arguments} to make an atomic formula.
\end{definition}

Terms and predicates are the only kinds of non-logical vocabulary that we will
make use of. Everything else is logic.

\begin{definition}[(Non-)Logical Vocabulary]
  The \emph{logical vocabulary} of first-order logic includes all of the
  brackets, quantifiers, operators and connectives. These do not vary between
  the languages of FOL. They all use the same \emph{syntax}.
  
  The \emph{non-logical vocabulary} of first-order logic includes all of the
  so-called \emph{content expressions} of the language: \emph{proper names} and
  \emph{predicates} of every arity. These will vary from language to language.

  Even though \emph{variables} and \emph{dummy names} can occur in the slots of
  predicates, they are \emph{logical vocabulary}. They are only ever used in the
  context of an associated quantifier.
\end{definition}

As with TFL, we can express generalizations about the logic using a consistent
metalogical notation. 

\begin{definition}[Metalogical Symbols]
  We will adopt the following conventions for the use of metalogical variables
  in describing first-order logic:
  \begin{argument}
  \item[\ealpha] for expressions (in general)
  \item[\etau] for terms (in general)
  \item[\enu] for proper names
  \item[\edelta] for dummy names
  \item[\exi] for variables
  \item[\ebeta] for predicates
  \item[\ephi] for well-formed formulas
  \end{argument}
\end{definition}

\subsection{Well-Formed Formulas}

Unlike TFL, FOL needs to define its atomic sentences since they are not mere
truth-bearers but have constituent parts.

\begin{definition}[Atomic Formula]
  \label{def:fol-atomic-formula}
  An \emph{atomic formula} of first-order logic is defined as:\\
  \begin{enumerate}
  \item Any zero-place predicate is an atomic formula.
  \item If \ebeta is an n-place predicate and $\etau_1\dots\etau_n$ are
    terms, then $\ebeta\etau_1,\dots,\etau_n$ is an atomic formula.
  \item If $\etau_1$ and $\etau_2$ are terms then $\etau_1=\etau_2$ is an
    atomic formula.
  \item Nothing else is an atomic formula. 
  \end{enumerate}
\end{definition}

As we will see next lecture, we can determine a truth value for everything that
we have defined as an atomic formula. Now that we have these basic
truth-bearers, we can fold in the components of truth-functional logic. All of
the TFL connectives will be included.

\begin{definition}[Well-Formed Formula]
  \label{def:fol-wff}
  A \emph{well-formed formula} (\emph{wff}) of first-order logic is defined
    recusively as:\\
  \begin{enumerate}
    \item Any atomic formula of FOL is a wff of FOL
    \item \ebot is a wff  of FOL.
    \item If \ephi is a wff of FOL, then \enot\ephi is a wff of FOL.
    \item If \ephi and \epsi are wffs of FOL, then $\ephi\eand\epsi$ is a
      wff of FOL.
    \item If \ephi and \epsi are wffs of FOL, then $\ephi\eor\epsi$ is a
      wff of FOL.
    \item If \ephi and \epsi are wffs of FOL, then $\ephi\eif\epsi$ is a
      wff of FOL.
    \item If \ephi and \epsi are wffs of FOL, then $\ephi\eiff\epsi$ is a
      wff of FOL.
    \item If \ephi(\etau) is a wff of FOL containing one or more occurrences of
      the term \etau and \ephi(\exi) is the expression which results from
      replacing \etau wherever it occurs in \ephi(\etau) with the new variable
      \exi, then \eexists\exi\ephi(\exi) and \eforall\exi\ephi(\exi) are wffs of
      FOL.
    \item Nothing else is a wff of FOL.
  \end{enumerate}
\end{definition}

We've added two new rules for quantified sentences. We will refer to
\eforall\exi and \eexists\exi as \emph{quantifiers}, to \eforall and \eexists
as \emph{quantifier symbols}, and to \exi as a \emph{variable}. \eforall\exi is
called \emph{universal} and \eexists is called \emph{existential}.

\subsection{Interlude: Free Variables vs Dummy Names}
\label{sec:fol-free-variables}

Notice that on this forumulation of the grammar, variables are artefacts of the
process that introduces quantifiers. If you've been following along in
\cite{Mag25}, you will notice that this departs from the rules presented there.
We will discuss the differences briefly before committing ourselves to the rules
in Def \ref{def:fol-wff}.

There are two important differences in the grammatical rules for formulas.
First, \cite{Mag25} allow variables to serve as \emph{terms}, making them
available in the construction of \emph{atomic formulas}. Second, they therefore
allow the construction of sentences which contain \emph{free variables}:
variables that are not within the scope of any quantifiers.

\begin{definition}[Well-Formed Formula (\cite{Mag25}]
  \label{def:fol-wff-mag25}
  A \emph{well-formed formula} (\emph{wff}) of first-order logic is defined
    recusively as:\\
    \begin{enumerate}
    \item Any atomic formula of FOL is a wff of FOL
    \item \ebot is a wff  of FOL.
    \item If \ephi is a wff of FOL, then \enot\ephi is a wff of FOL.
    \item If \ephi and \epsi are wffs of FOL, then $\ephi\eand\epsi$ is a
      wff of FOL.
    \item If \ephi and \epsi are wffs of FOL, then $\ephi\eor\epsi$ is a
      wff of FOL.
    \item If \ephi and \epsi are wffs of FOL, then $\ephi\eif\epsi$ is a
      wff of FOL.
    \item If \ephi and \epsi are wffs of FOL, then $\ephi\eiff\epsi$ is a
      wff of FOL.
    \item If \ephi is a wff or FOL and \exi is a variable, then
      \eforall\exi\ephi is a wff of FOL
    \item If \ephi is a wff or FOL and \exi is a variable, then
      \eexists\exi\ephi is a wff of FOL
    \item Nothing else is a wff of FOL.
    \end{enumerate}
\end{definition}

Notice how the quantifier rules do not require any particular relationship
between the variables in \ephi and the one added with the quantifier (\exi).
It's just that \eexists\exi will bind \exi if there happen to be any free in
\ephi. Maybe there are not. We cannot build any formulas containing unbound
variables using our rules of syntax. Variables are only ever introduced along
with a quantifier that binds them. That makes the following definition from
\cite{Mag25} more or less moot in our system. \textbf{All} of our variables are
\emph{bound} by quantifiers, from the moment they are introduced.

\begin{definition}[Bound and Free Variables]
  \label{def:fol-bound-variable}
  In first-order logic, a variable \exi is \emph{bound} iff it occurs in
  the scope of a quantifier that uses it: \eforall\exi or \eexists\exi.
  Otherwise it is \emph{free}.
\end{definition}

But note that we will make reference to expressions nominally containing free
variables when we talk about the \emph{scope} of a quantifier. For example, when
we talk about \ephi(\exi)--i.e. an expression in the scope of \eexists\exi in
\eexists\exi\ephi(\exi)--we note that it is an expression with \emph{free}
occurrences of \exi. But in the full context, \exi is bound by \eexists\exi. In
the rare instances that we talk about such expressions, we will refer to them
as \emph{open wffs}.

\subsection{Scope and Bound Variables}

Given the close alignment of the definitions, we can recycle some of our ideas
about operators and scope from TFL.

\begin{definition}[Main Logical Operator]
  \label{def:fol-main-operator}
  The main operator in a wff of first-order logic is the one that was
  introduced last, given the recursive application of the construction rules
  introduced in Def \ref{def:fol-wff}.
\end{definition}

\begin{definition}[Scope]
  \label{def:fol-scope}
  The \emph{scope} of an operator in first-order logic is the subformula for
  which it serves as the main logical operator.
\end{definition}

We can analyze these properties using the same toolkit as TFL, with some tweaks
to account for quantifiers. First, we can make use of parse trees to dissect
the structure of FOL wffs. We will need a new rule that allows us to add dummy
names wherever we encounter a quantified sentence.

\begin{definition}[Standard Parse Tree]
  Given a wff of first-order logic, we can build its \emph{standard parse tree}
  by using the following procedure:
  \begin{enumerate}
  \item If the current node is \enot\ephi add a node below it with \ephi
  \item If the current node is \ephi\eand\epsi, \ephi\eor\epsi,
    \ephi\eif\epsi, or \ephi\eiff\epsi, then split the tree with \ephi on
    one branch and \epsi on the other.
  \item If the current node is \eexists\exi\ephi(\exi) or
    \eforall\exi\ephi(\exi), then add a node below and with the instance
    \ephi($\edelta_i$) where $\edelta_i$ is the next available dummy name.
  \end{enumerate}
\end{definition}

Once again, the \emph{main logical operator} is the connective associated with
the first rule applied at the top of the tree (i.e. the last rule applied in the
construction of the wff). And the subformula that decorates the node above a
given rule application is the \emph{scope} of the associated operator. These
definitions work exaclty as they do in TFL, except for the addition of rules
for \eforall and \eexists, which are analogous to \enot in their scopal
behavior. In Def \ref{def:fol-formula}, the scope of both \eforall\exi and
\eexists\exi includes \ephi(\exi).

\subsection{Sentences and Dummy WFFs}

We can now identify sentences as a special kind of wff, where all the names are
proper and all the variables are bound. Well-formed formulas with dummy names
are called \emph{dummy wffs} and those without are called \emph{sentences}.

\begin{definition}[Sentence]
  \label{def:fol-sentence}
  In first-order logic, a wff \ephi is a \emph{sentence} iff all the terms that
  \ephi contains are proper names. Otherwise it is called a \emph{dummy wff}.
\end{definition}

We will adopt some bracketing simplification for sanity. For example,
$((A\eand B)\eand C)$ can be written $(A\eand B\eand C)$ and
$((A\eor B)\eor C)$ can be written $(A\eor B\eor C)$, since it will not
effect the meaning of the sentence to do so. We also allow the use of square
brackets ``['' and ``]'' as substitutes for parens ``('' and ``)'', to make
pairs of brackets easier to spot.

\subsection{Ambiguity and Quantifier Scope}

The quantifiers in first-order logic bind variables that occupy slots in the
structure of FOL atomic formulas that are typically reserved for terms. But
the structure of the syllogisms we reviewed last week is more like a comparison
of predicates: ``All F are G'' or ``Some F is G''. Our solution was to treat
the scope of the quantifiers as sentences connecting the two predicates we
care about, taking them to use the same variable, so:

\begin{argument}
  \item[ ] \eforall x(Fx\eif Gx)
  \item[ ] \eexists x(Fx\eand Gx)
\end{argument}

Indeed, by our rules, these sentences were constructed from \emph{instances}
containing dummy names:

\begin{argument}
  \item[ ] Fc\eif Gc
  \item[ ] Fa\eand Ga
\end{argument}

This approach gets interesting when we start to use multiple quantifiers and
predicates with a greater arity. We saw this informally when we tried to
symbolize some examples last lecture. ``No woman loves every man'' could be
symbolized in either of the following two ways, each with a different meaning.

\begin{argument}
\item[ ] \enot\eexists x(Gx\eand\eforall y(Fy\eif Lxy))
\item[ ] \eforall y(Fy \eif \enot\eexists x(Gx\eand Lxy))
\end{argument}

These formulas can be paraphrased, respectively, as:

\begin{argument}
\item[ ] (There is no woman such that)(Every man is such that)(She loves him)
\item[ ] (Every man is such that)(There is no woman such that)(She loves him)
\end{argument}

Despite all the syntactic apparatus, the difference really just comes down to
the order of the quantifiers: their scope. When the scope of a quantifer
$\ealpha_1$ contains another quantifier $\ealpha_2$ and its scope, we say that
$\ealpha_1$ takes \emph{wide scope} over $\ealpha_2$.

Now that we have a grammar to work with, we can see how FOL allows us to be
precise about some of the ambiguities of language that we just cannot
distinguish in a strictly sentence-level logic like TFL.
\cite[Chapter 29]{Mag25} works through some examples that are
particularly illustrative.

\begin{enumerate}
\item Everything that glitters is not gold.
\item Only young cats are playful.
\item Everyone went to see a movie.
\item The king of France is not bald.
\end{enumerate}

The first sentence has two readings, the basic \emph{every A is B} one and
another that more close reflects its idiomatic sense, something like
\emph{it is not the case that everything that glitters is gold}. 

\begin{enumerate}
\item $\eforall x(F(x)\eif\enot G(x))$
\item $\enot\eforall x(F(x)\eif G(x))$
\end{enumerate}

Note how the difference involves the relative scopes of negation and the
universal quantifier.

Similarly, for the second sentence, we get three distinct readings, normally
distinguished in ordinary English with patterns of word stress.

\begin{enumerate}
\item If an animal is playful, then it is a young cat.
\item If a cat is playful, then it is young.
\item If an animal is young anf playful, then it is a cat. 
\end{enumerate}

These can all be easily distinguished via symbolization in FOL, again by
taking account of the relative scopes of operators and quantifiers.

\begin{enumerate}
\item $\eforall x(P(x)\eif(Y(x)\eand C(x))$ 
\item $\eforall x((C(x)\eand P(x))\eif Y(x))$ 
\item $\eforall x((Y(x)\eand P(x))\eif C(x))$ 
\end{enumerate}

In the third sentence, we see how the quantifiers interact. 

\begin{enumerate}
\item There is a specific movie and everyone saw that movie.
\item Everyone went to some movie or other.
\end{enumerate}

The two readings differ with respect to which quantifier is placed within the
scope of the other.

\begin{enumerate}
\item $\eexists x(M(x)\eand\eforall y(P(y)\eif\S(y,x)))$
\item $\eforall y(P(y)\eif\forall x(M(x)\eand S(y,x)))$
\end{enumerate}

Finally, Russell's treatment of definite desriptions, while it avoids
\emph{truth-value gaps}, makes the fourth sentence ambiguous in a way that
directly impacts its truth value.

\begin{enumerate}
\item There is a king of France and he is not bald. (False)
\item There is no king of France who is bald. (True)
\end{enumerate}

And, once again, the differences come down to which operator or quantifier
takes \emph{wide scope}. 

\begin{enumerate}
\item $\eexists x(K(x)\eand\forall y(K(y)\eif x\neq y)\eand\enot B(x))$
\item $\enot\eexists x(K(x)\eand\forall y(K(y)\eif x\neq y)\eand B(x))$
\end{enumerate}

% more scope examples for definites (pp378-81)

Despite all these structural ambiguities, there are some cases where one
quantifier taking wide scope over another is not meaningfully different than
the opposite. First, we can always reverse the position of two adjacent
quantifiers of the same type.

\begin{observation}[Adjacent Quantifiers]
  \label{obs:adjacent-quantifiers}
  If two adjacent quantifiers are the same kind (either both universal or both
  existential), then it does not matter what order they occur in. In other
  words, $\eforall\exi_1\eforall\exi_2\ephi$ is equivalent to
  $\eforall\exi_2\eforall\exi_1\ephi$ and
  $\eexists\exi_1\eexists\exi_2\ephi$ is equivalent to
  $\eexists\exi_2\eexists\exi_1\ephi$ for all $\ephi$.
\end{observation}

Similarly, we can adjust the scope of a quantifier in many cases where its
variable does not occur in the subformulas impacted by the change.

\begin{observation}[Moving Quantifiers]
  \label{obs:moving-quantifiers}
  When \epsi does not contain the variable \exi, the following pairs of wffs
  are equivalent.
  \begin{enumerate}
  \item \eforall\exi\ephi(\exi)\eand\epsi and
    \eforall\exi(\ephi(\exi)\eand\epsi)
  \item \eexists\exi\ephi(\exi)\eand\epsi and
    \eexists\exi(\ephi(\exi)\eand\epsi)
  \item \eforall\exi\ephi(\exi)\eor\epsi and
    \eforall\exi(\ephi(\exi)\eor\epsi)
  \item \eexists\exi\ephi(\exi)\eor\epsi and
    \eexists\exi(\ephi(\exi)\eor\epsi)
  \item \eforall\exi\ephi(\exi)\eif\epsi and
    \eexists\exi(\ephi(\exi)\eif\epsi)
  \item \eexists\exi\ephi(\exi)\eif\epsi and
    \eforall\exi(\ephi(\exi)\eif\epsi)
  \item \epsi\eif\eforall\exi\ephi(\exi) and
    \eforall\exi(\epsi\eif\ephi(\exi))
  \item \epsi\eif\eexists\exi\ephi(\exi) and
    \eexists\exi(\epsi\eif\ephi(\exi))
  \end{enumerate}
\end{observation}

%% adding identity (p367)
Now that we know how to make them, next week we will discuss how to interpret
sentences of first-order logic.
